%%=============================================================================
%% Conclusie
%%=============================================================================

\chapter{Conclusie \& Discussie}%
\label{ch:conclusie_discussie}

% TODO: Trek een duidelijke conclusie, in de vorm van een antwoord op de
% onderzoeksvra(a)g(en). Wat was jouw bijdrage aan het onderzoeksdomein en
% hoe biedt dit meerwaarde aan het vakgebied/doelgroep? 
% Reflecteer kritisch over het resultaat. In Engelse teksten wordt deze sectie
% ``Discussion'' genoemd. Had je deze uitkomst verwacht? Zijn er zaken die nog
% niet duidelijk zijn?
% Heeft het onderzoek geleid tot nieuwe vragen die uitnodigen tot verder 
%onderzoek?

% \lipsum[76-80]

\section{Conclusie}%
\label{sec:conclusie}

Om op de centrale onderzoeksvraag van deze bachelorproef te antwoorden, is een proof-of-concept converter ontwikkeld genaamd PlantUML2Ansible, die netwerk- en implementatiediagrammen, geschreven in PlantUML-code, als invoer aanneemt en een volledig IaC- en `\emph{configuration management}'-ondersteunde omgeving (aangeboden door Vagrant en Ansible, respectievelijk) als uitvoer produceert.

Naast dit primaire gebruiksscenario (\nameref{subsec:fullmode}) is er ook een alternatief scenario ontwikkeld om enkel een Vagrant-omgeving te produceren (\nameref{subsec:iaconlymode}).

Bestaande labo-omgevingen uit bepaalde OLODS binnen de opleiding \emph{Toegepaste Informatica} (\emph{Cybersecurity Advanced} en \emph{Infrastructure Automation}) zijn herwerkt als bruikbare invoer voor de converter (zijnde netwerk- en/of implementatiediagrammen) en zijn, binnen de beperkingen van het script, met succes omgezet naar functionele omgevingen.
De initiële labo-omgeving van \emph{Cybersecurity Advanced} is via de ``IaC-only mode'' volledig gereproduceerd kunnen worden als  een Vagrant-ondersteunde omgeving.
De labo's 2 en 3 uit \emph{Infrastructure Automation} zijn ook via de ``full mode'' gerecreëerd als Vagrant- en Ansible-ondersteunde omgeving, maar zonder de router \texttt{r001}.

Deze resultaten tonen dus aan dat het technisch haalbaar is om, op basis van PlantUML-bronbestanden, volledige en functionele omgevingen te genereren -- zonder manuele tussenkomst van de gebruiker.

Tegelijkertijd kent de proof-of-concept een aantal beperkingen die de scope van het huidige onderzoek afbakenen. De ``full mode'' ondersteunt slechts één netwerk, het ondersteunde gastbesturingssysteem wordt beperkt tot \emph{AlmaLinux} 9, en de invoerdiagrammen moeten voldoen aan een welomschreven en beperkt subset van de PlantUML-syntax.

De meerwaarde van deze oplossing situeert zich voornamelijk in de educatieve context waarvoor ze ontworpen is. Docenten binnen de opleiding \emph{Toegepaste Informatica} op HOGENT kunnen met deze tool sneller nieuwe labo-omgevingen ontwerpen en deze tegelijk voorzien van functionele code om deze effectief op te zetten en te configureren. De kloof tussen de ontwerpfase (UML-diagram) en de implementatiefase (uitgerolde omgeving) wordt zo met andere woorden aanzienlijk verkleind.

Voor studenten kan de tool de link tussen deze twee fases juist alsmaar meer verduidelijken, en de drempel verlagen om nieuwe infrastructuur-gerelateerde ideeën al experimenterend vorm te geven.

\section{Discussie}%
\label{sec:discussie}

De beperkingen van deze proof-of-concept leggen verschillende fundamenten waarop toekomstige uitbreidingen verder gebouwd kunnen worden.

Het ondersteunen voor meerdere netwerken in beide modi is technisch mogelijk, mits de configuratie van routering (en de gerelateerde firewallregels) wordt verwerkt in de logica van het script. Daarnaast moet er, onder andere, een uitbreiding van de DNS-zone-generatielogica uitgevoerd worden, alsook een herziene aanpak voor de toewijzing (van het IP-adres) van de Ansible \texttt{control}-node.

De beperking van \emph{AlmaLinux} 9 als gastbesturingssysteem kan geremedieerd worden door de configuratie van de rollen niet meer afhankelijk te maken van de bestaande rollen die gebruikt worden binnen de opleiding \emph{Toegepaste Informatica}. Door over te stappen naar meer distributie-agnostische modules in deze  nieuwe rollen, zou de converter ook andere distributies kunnen ondersteunen, waaronder de nieuwere \emph{AlmaLinux} 10 en misschien zelfs Ubuntu Server of Alpine Linux.

Het huidige \texttt{role-config.yml}-bestand bevat momenteel slechts een beperkt aantal voorgedefinieerde rollen. Hoewel het bestand alvast een modulaire basis heeft, is het op dit moment niet mogelijk om nieuwe rollen (zelfs met de minste complexiteit) te definiëren zonder de broncode van het script te wijzigen.

Andere, minder cruciale uitbreidingen zouden de volgende kunnen inhouden:
\begin{itemize}
    \item Compatibiliteit met IPv6
    \item Ondersteuning voor andere IaC- (Terraform, OpenTofu, \ldots) en `\emph{configuration management}'-backends (Puppet, Chef, \ldots)
    % \item Grafische gebruikersinterface of integraties met IDE's via plugins (bijvoorbeeld een `\emph{Visual Studio Code}'-extensie)
    \item Formele validatie van de PlantUML-invoer en de Vagrant- en Ansible-uitvoer via geïntegreerde linters
\end{itemize}